\chapter{Conclusão}
\label{cap:conclusao}
%\phantom{2}

As altas taxas de incidência do câncer de rim em todo o mundo mostram a importância do desenvolvimento de pesquisas que forneçam suporte ao diagnóstico precoce da doença. Essa é uma etapa fundamental para viabilizar o tratamento adequado aos pacientes. Nesse contexto, este trabalho apresentou um método automático para a tarefa de segmentação de rins e tumores renais baseado em aprendizado profundo e técnicas de processamento de imagens para reduzir os falsos positivos.



\section{Contribuições}
\label{sec:contribuições}

Finalmente, as principais contribuições do método proposto desenvolvido nesta tese são descritas a seguir:

\begin{enumerate}
    \item Desenvolvimento de um método automático para agrupar casos (exames) e distribuí-los proporcionalmente no conjunto de treinamento e validação para garantir que os modelos fossem equilibrados e obtivessem melhor desempenho;
    
    \item Adaptação da arquitetura DeepLabv3+, com o uso da DPN-31 como codificador para a segmentação de tumores renais, fornecendo resultados precisos;
    
   
\end{enumerate}

\section{Trabalhos Futuros}
\label{sec:trabalhos-futuros}

Apesar dos bons resultados obtidos para a segmentação de rins e tumores renais em imagens de TCs, melhorias ainda podem ser feitas no método proposto visando incrementar a sua eficiência. A seguir são destacadas algumas sugestões:

\begin{enumerate}
    \item Considerando a diversidade de aquisições da base de imagens, padronizar as entradas da rede usando uma rede \textit{autoencoder} como um pré-processamento poderia aumentar a eficiência dos modelos de segmentação;
    
    \item Assim como na abordagem 2.5D usada no método proposto as fatias são visualizadas sequencialmente. Acredita-se que o uso de recursos recorrentes da rede neural pode melhorar o desempenho, pois essas redes podem “lembrar” informações de fatias anteriores, adicionando recursos ao modelo;
    
  
\end{enumerate}

\section{Produções Científicas}
\label{sec:producoes-cientificas}

A Tabela~\ref{tab:artigos-autoria} apresenta os artigos diretamente relacionados ao método proposto ....... Além disso, a Tabela~\ref{tab:artigos-coautoria} lista os artigos científicos publicados e submetidos em outras aplicações de processamento de imagens e visão computacional desde o início do doutorado.

\begin{table}[ht!]
\centering
\onehalfspacing
\caption{Produções científicas em relação ao método proposto para segmentação de rins e tumores renais.}
\label{tab:artigos-autoria}
\resizebox{\columnwidth}{!}{
\begin{tabular}{p{10cm}lccc}
\hline
\textbf{Artigo}                                                                                                                         & \centering \textbf{Tipo} & \textbf{Qualis} & \textbf{Status} \\ \hline
Titulo do artigo Em: Nome do periodico Ano: 2020.        & Periódico     & A1              & Publicado       \\ \hline
Titulo do artigo Em: Nome do periodico Ano: 2020.      & Periódico     & A1              & Publicado      \\ \hline
\end{tabular}
}
\end{table}

\begin{table}[ht!]
\centering
\onehalfspacing
\caption{Produções científicas em outras aplicações de processamento de imagens e visão computacional.}
\label{tab:artigos-coautoria}
\resizebox{\columnwidth}{!}{
\begin{tabular}{p{11cm}lccc}
\hline
\textbf{Artigo}                                                                                                                                                                            & \textbf{Tipo} & \textbf{Qualis} & \textbf{Status} \\ \hline
Titulo do artigo Em: Nome do periodico Ano: 2020.      . & Periódico     & A1              & Publicado       \\ \hline
Titulo do artigo Em: Nome do periodico Ano: 2020.                          & Simpósio      & B3              & Publicado       \\ \hline
\end{tabular}
}
\end{table}